Para el objetivo de la identificación y separación de elementos en una obra musical, se ha utilizado una arquitectura \textit{LSTM}, en particular la de \textit{Open-Unminx} \cite{rafii18}. Su artículo de presentación incluye una sección didáctica donde se ha aprendido gran parte del método usado. 

En el desarrollo se han contemplado dos escenarios de separación:

\begin{itemize}
    \item 2 fuentes: voz/acompañamiento
    \item 4 fuentes: voz/bajo/batería/otros
\end{itemize}

En dos fuentes, cuando la entrada al proceso es de una canción con mas de dos elementos , todo lo que no sea voz se agrupa en el acompañamiento.
